%------------------------------------------------------------------------------%
\documentclass[fleqn,utf8,aspectratio=169,12pt]{beamer}

\usepackage{gaal}

\title{Matrizes Inversas}

%------------------------------------------------------------------------------%
\begin{document}

\addtitle

%------------------------------------------------------------------------------%
\section{Matrizes Inversas}

%------------------------------------------------------------------------------%
\begin{frame}{Matriz inversa}
  Dada uma matriz quadrada $A$

  \pause
  $A$ é \emph{invertível}
  \pause
  quando existe uma matriz
  \emph{$A^{-1}$}

  \pause
  tal que
  \[
    \pause
    A A^{-1} = I
  \]
  \[
    \pause
    A^{-1} A = I
  \]
  \pause
  $I$ é a matriz identidade
\end{frame}

%------------------------------------------------------------------------------%
%------------------------------------------------------------------------------%
\begin{question}
  Dadas as matrizes
  \[
    A = \begin{pmatrix}
      1 & 2 & 0 \\
      0 & 1 & 1 \\
      1 & 0 & 1
    \end{pmatrix}
    \quad
    \text{e}
    \quad
    B = \begin{pmatrix}
      1  & -2 & 2  \\
      1  & 1  & -1 \\
      -1 & 2  & 1
    \end{pmatrix}
  \]
  verifique se \(B\) é a matriz inversa de \(A\)
\end{question}

%------------------------------------------------------------------------------%
\begin{resolution}{Solução}
  Precisamos
  \begin{enumerate}
      \item Calcular o produto \(AB\)
      \item Calcular o produto \(BA\)
      \item Comparar os resultados obtidos com a matriz identidade \(I_3\)
  \end{enumerate}
\end{resolution}

%------------------------------------------------------------------------------%
\begin{resolution}{{Cálculo de \(AB\)}}
  \begin{align*}
    \onslide<+->{AB}
    \onslide<+->{& =
    \begin{pmatrix}
      1 & 2 & 0 \\
      0 & 1 & 1 \\
      1 & 0 & 1
    \end{pmatrix}
    \;
    \begin{pmatrix}
      1  & -2 & 2  \\
      1  & 1  & -1 \\
      -1 & 2  & 1
    \end{pmatrix}}
    \onslide<+->{\\ & =
    \begin{pmatrix}
      1\cdot1+2\cdot1+0\cdot(-1) & 1\cdot(-2)+2\cdot1+0\cdot2 & 1\cdot2+2\cdot(-1)+0\cdot1 \\
      0\cdot1+1\cdot1+1\cdot(-1) & 0\cdot(-2)+1\cdot1+1\cdot2 & 0\cdot2+1\cdot(-1)+1\cdot1 \\
      1\cdot1+0\cdot1+1\cdot(-1) & 1\cdot(-2)+0\cdot1+1\cdot2 & 1\cdot2+0\cdot(-1)+1\cdot1
    \end{pmatrix}}
    \onslide<+->{\\ & =
    \begin{pmatrix}
      1+2+0 & -2+2+0 & 2-2+0 \\
      0+1-1 & 0+1+2  & 0-1+1 \\
      1+0-1 & -2+0+2 & 2+0+1
    \end{pmatrix}}
  \end{align*}
\end{resolution}

%------------------------------------------------------------------------------%
\begin{resolution}{{Cálculo de \(AB\)}}
  \begin{align*}
    \onslide<+->{AB & =
    \begin{pmatrix}
      1+2+0 & -2+2+0 & 2-2+0 \\
      0+1-1 & 0+1+2  & 0-1+1 \\
      1+0-1 & -2+0+2 & 2+0+1
    \end{pmatrix}}
    \onslide<+->{\\ & =
    \begin{pmatrix}
      3 & 0 & 0 \\
      0 & 3 & 0 \\
      0 & 0 & 3
    \end{pmatrix}}
    \onslide<+->{\\ & =
    3I_3}
  \end{align*}
\end{resolution}

%------------------------------------------------------------------------------%
\begin{resolution}{{Cálculo de \(BA\)}}
  \begin{align*}
    \onslide<+->{BA}
    \onslide<+->{& =
    \begin{pmatrix}
      1  & -2 & 2  \\
      1  & 1  & -1 \\
      -1 & 2  & 1
    \end{pmatrix}
    \;
    \begin{pmatrix}
      1 & 2 & 0 \\
      0 & 1 & 1 \\
      1 & 0 & 1
    \end{pmatrix}}
    \onslide<+->{\\ & =
    \begin{pmatrix}
      1\cdot1+(-2)\cdot0+2\cdot1 & 1\cdot2+(-2)\cdot1+2\cdot0 & 1\cdot0+(-2)\cdot1+2\cdot1 \\
      1\cdot1+1\cdot0+(-1)\cdot1 & 1\cdot2+1\cdot1+(-1)\cdot0 & 1\cdot0+1\cdot1+(-1)\cdot1 \\
      (-1)\cdot1+2\cdot0+1\cdot1 & (-1)\cdot2+2\cdot1+1\cdot0 & (-1)\cdot0+2\cdot1+1\cdot1
    \end{pmatrix}}
    \onslide<+->{\\ & =
    \begin{pmatrix}
      1+0+2  & 2-2+0  & 0-2+2 \\
      1+0-1  & 2+1+0  & 0+1-1 \\
      -1+0+1 & -2+2+0 & 0+2+1
    \end{pmatrix}}
  \end{align*}
\end{resolution}


%------------------------------------------------------------------------------%
\begin{resolution}{{Cálculo de \(BA\)}}
  \begin{align*}
    \onslide<+->{BA & =
    \begin{pmatrix}
      1+0+2  & 2-2+0  & 0-2+2 \\
      1+0-1  & 2+1+0  & 0+1-1 \\
      -1+0+1 & -2+2+0 & 0+2+1
    \end{pmatrix}}
    \onslide<+->{\\ & =
    \begin{pmatrix}
      3 & 0 & 0 \\
      0 & 3 & 0 \\
      0 & 0 & 3
    \end{pmatrix}}
    \onslide<+->{\\ & =
    3I_3}
  \end{align*}
\end{resolution}

%------------------------------------------------------------------------------%
\begin{resolution}{Solução}
  Como
  \[
    AB = 3I_3 \neq I_3
  \]
  \[
    BA = 3I_3 \neq I_3
  \]
  \pause
  $B$ não é a inversa de $A$
\end{resolution}

%------------------------------------------------------------------------------%
\begin{resolution}{Solução}
  Porém
  \[
    \pause
    A^{-1} = \frac{1}{3} B
  \]
  \pause
  pois
  \[
    \pause
    AA^{-1}
    \pause
    \;=\;
    A \frac{1}{3} B
    \pause
    \;=\;
    \frac{1}{3} AB
    \pause
    \;=\;
    \frac{1}{3} \; 3 \; I_3
    \pause
    \;=\;
    I_3
  \]
  \[
    \pause
    A^{-1}A
    \pause
    \;=\;
    \frac{1}{3} BA
    \pause
    \;=\;
    \frac{1}{3} \; 3 \; I_3
    \pause
    \;=\;
    I_3
  \]
\end{resolution}

%------------------------------------------------------------------------------%
\endinput
\subsection*{Comparação com a identidade}


\subsection*{Conclusão}

Como
\[
A \cdot B = B \cdot A = 3I_3 \neq I_3,
\]
concluímos que \(B\) \textbf{não} é a matriz inversa de \(A\).

A inversa de \(A\), caso exista, seria \(\dfrac{1}{3}B\).

\end{document}

%------------------------------------------------------------------------------%
\section{Calculando a Matriz Inversa}

%------------------------------------------------------------------------------%
\begin{frame}{Calculando a Matriz Inversa}
  Podemos calcular a matriz inversa usando o escalonamento de Gauss-Jordan

  \pause
  Começamos com a matriz aumentada com a identidade

  \[
    \left[ \;A\; \middle| \;I\; \right]
  \]

  \pause
  Escalonamos $A$

  \pause
  O método transforma $I$ em $A^{-1}$
\end{frame}

%------------------------------------------------------------------------------%
%------------------------------------------------------------------------------%
\begin{question}
  Encontre a inversa de
  \quad
  \(
    A =
    \begin{bmatrix}
      2 & 1 \\
      1 & 1
    \end{bmatrix}
  \)
\end{question}

%------------------------------------------------------------------------------%
\begin{resolution}{Matriz Aumentada}
  Matriz aumentada \([\,A\,|\,I\,]\)
  \pause
  \[
    \begin{amatrix}[2]{2}
      \emph{2} & \emph{1} & 1 & 0 \\
      \emph{1} & \emph{1} & 0 & 1
    \end{amatrix}
  \]

  \pause
  Queremos transformar a parte da esquerda na identidade $I_2$
\end{resolution}

%------------------------------------------------------------------------------%
\begin{resolution}{Escalonamento}
\begin{columns}
\column{0.3\linewidth}
  \[\;\;\;
    \begin{amatrix}[2]{2}
      2 & 1 & 1 & 0 \\
      1 & 1 & 0 & 1
    \end{amatrix}
  \]
  \pause
\column{0.7\linewidth}
  Trocamos as linhas
  \\ \pause
  \(
    L_1\leftrightarrow L_2
  \)
\end{columns}

\pause
\[
  \begin{amatrix}[2]{2}
    1 & 1 & 0 & 1 \\
    2 & 1 & 1 & 0
  \end{amatrix}
\]
\end{resolution}

%------------------------------------------------------------------------------%
\begin{resolution}{Escalonamento}
\begin{columns}
\column{0.3\linewidth}
  \[\;\;\;
    \begin{amatrix}[2]{2}
      1 & 1 & 0 & 1 \\
      2 & 1 & 1 & 0
    \end{amatrix}
  \]
  \pause
\column{0.7\linewidth}
  Eliminamos os elementos abaixo do primeiro pivô
  \\ \pause
  \(
    L_2 \leftarrow L_2 - 2 L_1
  \)
\end{columns}
\vspace{-\baselineskip}
\[
  \pause
  L'_2
  \;=\;
  L_2 - 2 L_1
  \pause
  \;=\;
  \begin{amatrix}[2]{2}
    2 & 1 & 1 & 0
  \end{amatrix}
  - 2
  \begin{amatrix}[2]{2}
    1 & 1 & 0 & 1
  \end{amatrix}
  \pause
  \;=\;
  \begin{amatrix}[2]{2}
    0 & -1 & 1 & -2
  \end{amatrix}
\]
\pause
\[
  \begin{amatrix}[2]{2}
    1 & 1  & 0 & 1  \\
    0 & -1 & 1 & -2
  \end{amatrix}
\]
\end{resolution}

%------------------------------------------------------------------------------%
\begin{resolution}{Escalonamento}
\begin{columns}
\column{0.3\linewidth}
  \[\;\;\;
    \begin{amatrix}[2]{2}
      1 & 1  & 0 & 1  \\
      0 & -1 & 1 & -2
    \end{amatrix}
  \]
  \pause
\column{0.7\linewidth}
  Tornando o segundo pivô igual a 1
  \\ \pause
  \(
    L_2 \leftarrow - L_2
  \)
\end{columns}
\vspace{-\baselineskip}
\[
  \pause
  L'_2
  \;=\;
   - L_2
  \pause
  \;=\;
  -
  \begin{amatrix}[2]{2}
    0 & -1 & 1 & -2
  \end{amatrix}
  \pause
  \;=\;
  \begin{amatrix}[2]{2}
    0 & 1 & -1 & 2
  \end{amatrix}
\]
\pause
\[
  \begin{amatrix}[2]{2}
    1 & 1 & 0  & 1 \\
    0 & 1 & -1 & 2
  \end{amatrix}
\]
\end{resolution}

%------------------------------------------------------------------------------%
\begin{resolution}{Escalonamento}
\begin{columns}
\column{0.3\linewidth}
  \[\;\;\;
    \begin{amatrix}[2]{2}
      1 & 1 & 0  & 1 \\
      0 & 1 & -1 & 2
    \end{amatrix}
  \]
  \pause
\column{0.7\linewidth}
  Eliminamos os elementos acima do segundo pivô
  \\ \pause
  \(
    L_1 \leftarrow L_1 - L_2
  \)
\end{columns}
\vspace{-\baselineskip}
\[
  \pause
  L'_1
  \;=\;
  L_1 - L_2
  \pause
  \;=\;
  \begin{amatrix}[2]{2}
    1 & 1 & 0  & 1
  \end{amatrix}
  -
  \begin{amatrix}[2]{2}
    0 & 1 & -1 & 2
  \end{amatrix}
  \pause
  \;=\;
  \begin{amatrix}[2]{2}
    1 & 0 & 1  & -1
  \end{amatrix}
\]
\pause
\[
  \begin{amatrix}[2]{2}
    1 & 0 & 1  & -1 \\
    0 & 1 & -1 & 2
  \end{amatrix}
\]
\end{resolution}

%------------------------------------------------------------------------------%
\begin{resolution}{Solução}
  \[
    \begin{amatrix}[2]{2}
      1 & 0 & 1  & -1 \\
      0 & 1 & -1 & 2
    \end{amatrix}
  \]

  \pause
  Como a matriz tem posto completo

  \pause
  a inversa exite e pode ser lida do lado direito
  \pause
  \[
    A^{-1} =
    \begin{bmatrix}
      1  & -1 \\
      -1 & 2
    \end{bmatrix}
  \]
\end{resolution}

%------------------------------------------------------------------------------%
\begin{resolution}{Verificação}
  \onslide<+->{Podemos verificar o resultado calculando}
  \begin{align*}
    \onslide<+->{AA^{-1}}
    \onslide<+->{& =
      \begin{bmatrix}
        2 & 1 \\
        1 & 1
      \end{bmatrix}
      \begin{bmatrix}
        1  & -1 \\
        -1 & 2
      \end{bmatrix}}
    \\
    \onslide<+->{& =
      \begin{bmatrix}
        2\times 1 + 1 (-1) &  2(-1) + 1 \times 2\\
        1\times 1 + 1 (-1) &  1(-1) + 1 \times 2
      \end{bmatrix}}
    \\
    \onslide<+->{& =
      \begin{bmatrix}
        1 & 0 \\
        0 & 1
      \end{bmatrix}}
      \\
    \onslide<+->{& = I}
  \end{align*}
  \onslide<+->{Portanto, a matriz encontrada é realmente a inversa de $A$}
\end{resolution}

%------------------------------------------------------------------------------%
\endinput

%------------------------------------------------------------------------------%
\begin{question}
  Encontre a inversa da matriz
  \quad
  \(
    A =
    \begin{pmatrix}
      1 & 1 & 1 \\
      0 & 1 & 1 \\
      0 & 0 & 1
    \end{pmatrix}
  \)
\end{question}

%------------------------------------------------------------------------------%
\begin{resolution}{Matriz Aumentada}
  Matriz aumentada
  \[
    \begin{amatrix}[3]{3}
      1 & 1 & 1 & 1 & 0 & 0 \\
      0 & 1 & 1 & 0 & 1 & 0 \\
      0 & 0 & 1 & 0 & 0 & 1
    \end{amatrix}
  \]
\end{resolution}

%------------------------------------------------------------------------------%
\begin{resolution}{Escalonamento}
\begin{columns}
\column{0.3\linewidth}
  \[\;\;\;
    \begin{amatrix}[3]{3}
      1 & 1 & 1 & 1 & 0 & 0 \\
      0 & 1 & 1 & 0 & 1 & 0 \\
      0 & 0 & 1 & 0 & 0 & 1
    \end{amatrix}
  \]
  \pause
\column{0.7\linewidth}
  Eliminamos os elementos acima do terceiro pivô
  \\ \pause
  \(
    L_2 \leftarrow L_2 - L_3
  \)
  \\ \pause
  \(
    L_1 \leftarrow L_1 - L_3
  \)
\end{columns}
\vspace{-\baselineskip}
\[
  \hspace{-3ex}
  \pause
  L'_2
  \;=\;
  L_2 - L_3
  \pause
  \;=\;
  \begin{amatrix}[3]{3}
    0 & 1 & 1 & 0 & 1 & 0
  \end{amatrix}
  -
  \begin{amatrix}[3]{3}
    0 & 0 & 1 & 0 & 0 & 1
  \end{amatrix}
  \pause
  \;=\;
  \begin{amatrix}[3]{3}
    0 & 1 & 0 & 0 & 1 & -1
  \end{amatrix}
\]
\[
  \hspace{-3ex}
  \pause
  L'_1
  \;=\;
  L_1 - L_3
  \pause
  \;=\;
  \begin{amatrix}[3]{3}
    1 & 1 & 1 & 1 & 0 & 0
  \end{amatrix}
  -
  \begin{amatrix}[3]{3}
    0 & 0 & 1 & 0 & 0 & 1
  \end{amatrix}
  \pause
  \;=\;
  \begin{amatrix}[3]{3}
    1 & 1 & 0 & 1 & 0 & -1
  \end{amatrix}
\]
\pause
\[
  \begin{amatrix}[3]{3}
    1 & 1 & 0 & 1 & 0 & -1 \\
    0 & 1 & 0 & 0 & 1 & -1 \\
    0 & 0 & 1 & 0 & 0 & 1
  \end{amatrix}
\]
\end{resolution}

%------------------------------------------------------------------------------%
\begin{resolution}{Escalonamento}
\begin{columns}
\column{0.3\linewidth}
  \[\;\;\;
    \begin{amatrix}[3]{3}
      1 & 1 & 0 & 1 & 0 & -1 \\
      0 & 1 & 0 & 0 & 1 & -1 \\
      0 & 0 & 1 & 0 & 0 & 1
    \end{amatrix}
  \]
  \pause
\column{0.7\linewidth}
  Eliminamos os elementos acima do terceiro pivô
  \\ \pause
  \(
    L_1 \leftarrow L_1 - L_2
  \)
\end{columns}
\vspace{-\baselineskip}
\[
  \pause
  L'_1
  \;=\;
  L_1 - L_2
  \pause
  \;=\;
  \begin{amatrix}[3]{3}
    1 & 1 & 0 & 1 & 0 & -1
  \end{amatrix}
  -
  \begin{amatrix}[3]{3}
    0 & 1 & 0 & 0 & 1 & -1
  \end{amatrix}
  \pause
  \;=\;
  \begin{amatrix}[3]{3}
    1 & 0 & 0 & 1 & -1 & 0
  \end{amatrix}
\]
\pause
\[
  \begin{amatrix}[3]{3}
    1 & 0 & 0 & 1 & -1 & 0  \\
    0 & 1 & 0 & 0 & 1  & -1 \\
    0 & 0 & 1 & 0 & 0  & 1
  \end{amatrix}
\]
\end{resolution}

%------------------------------------------------------------------------------%
\begin{resolution}{Solução}
  \[
    \begin{amatrix}[3]{3}
      1 & 0 & 0 & 1 & -1 & 0  \\
      0 & 1 & 0 & 0 & 1  & -1 \\
      0 & 0 & 1 & 0 & 0  & 1
    \end{amatrix}
  \]

  \pause
  Como a matriz tem posto completo

  \pause
  A inversa existe e pode ser lida do lado direito
  \pause
  \[
    A^{-1} =
    \begin{bmatrix}
      1 & -1 & 0  \\
      0 & 1  & -1 \\
      0 & 0  & 1
    \end{bmatrix}
  \]
\end{resolution}

%------------------------------------------------------------------------------%
\endinput

%------------------------------------------------------------------------------%
\begin{question}
  Tente calcular a inversa da matriz
  \quad
  \(
    A =
    \begin{pmatrix}
      1 & 2 \\
      2 & 4
    \end{pmatrix}
  \)
\end{question}

%------------------------------------------------------------------------------%
\begin{resolution}{Matriz Aumentada}
  Matriz aumentada \([\,A\,|\,I\,]\)
  \pause
  \[
    \begin{amatrix}[2]{2}
      1 & 2 & 1 & 0 \\
      2 & 4 & 0 & 1
    \end{amatrix}
  \]

  \pause
  Queremos transformar a parte da esquerda na identidade $I_2$
\end{resolution}

%------------------------------------------------------------------------------%
\begin{resolution}{Escalonamento}
\begin{columns}
\column{0.3\linewidth}
  \[\;\;\;
    \begin{amatrix}[2]{2}
      1 & 2 & 1 & 0 \\
      2 & 4 & 0 & 1
    \end{amatrix}
  \]
  \pause
\column{0.7\linewidth}
  Eliminamos os elementos abaixo do primeiro pivô
  \\ \pause
  \(
    L_2 \leftarrow L_2 - 2 L_1
  \)
\end{columns}
\vspace{-\baselineskip}
\[
  \pause
  L'_2
  \;=\;
  L_2 - 2 L_1
  \pause
  \;=\;
  \begin{amatrix}[2]{2}
    2 & 4 & 0 & 1
  \end{amatrix}
  - 2
  \begin{amatrix}[2]{2}
    1 & 2 & 1 & 0
  \end{amatrix}
  \pause
  \;=\;
  \begin{amatrix}[2]{2}
    0 & 0 & -2 & 1
  \end{amatrix}
\]
\pause
\[
  \begin{amatrix}[2]{2}
    1 & 2 & 1 & 0 \\
    0 & 0 & -2 & 1
  \end{amatrix}
\]
\end{resolution}

%------------------------------------------------------------------------------%
\begin{resolution}{Solução}
  \[
    \begin{amatrix}[2]{2}
      1 & 2 & 1 & 0 \\
      0 & 0 & -2 & 1
    \end{amatrix}
  \]

  \pause
  O escalonamento da matriz $A$ produziu uma linha nula

  \pause
  A matriz $A$ não tem posto completo

  \pause
  A matriz $A$ é singular

  \pause
  A matriz $A$ não possui inversa
\end{resolution}

%------------------------------------------------------------------------------%
\endinput


%------------------------------------------------------------------------------%
\begin{frame}{Conexão Importante}
  \onslide<+->{Para uma matriz quadrada $A_{n}$ é equivalente}
  \begin{itemize}[<+->]
  \setlength{\itemsep}{1em}
  \item
    \(
      A \text{ é invertível}
    \)
  \item
    \(
      A \mathbf{x} = \mnula
    \)
    possui somente a solução trivial
  \item
    \(
      \posto(A) = n
    \)
  \end{itemize}
\end{frame}

%------------------------------------------------------------------------------%
\begin{frame}{Inversa e Sistemas Lineares}
  \onslide<+->{Dado
  \[
    A \mathbf{x} = \mathbf{b}
  \]}
  \onslide<+->{Se $A$ é invertível,}
  \onslide<+->{podemos multiplicar os dois lados por $A^{-1}$}
  \begin{align*}
    \onslide<+->{A \mathbf{x}        & = \mathbf{b}        \\}
    \onslide<+->{A^{-1} A \mathbf{x} & = A^{-1} \mathbf{b} \\}
    \onslide<+->{I \mathbf{x}        & = A^{-1} \mathbf{b} \\}
    \onslide<+->{\mathbf{x}          & = A^{-1} \mathbf{b}}
  \end{align*}
\end{frame}

%------------------------------------------------------------------------------%
%------------------------------------------------------------------------------%
\begin{question}
  Resolva o sistema linear utilizando a inversa
  \[
    \systeme{
      2x + y = 5,
       x + y = 3
    }
  \]
\end{question}

%------------------------------------------------------------------------------%
\begin{resolution}{Forma Matricial}
  Na forma matricial
  \[
    \begin{pmatrix}
      2 & 1 \\
      1 & 1
    \end{pmatrix}
    \begin{pmatrix} x \\ y \end{pmatrix} =
    \begin{pmatrix} 5 \\ 3 \end{pmatrix}
  \]

  \pause
  Equivalentemente \(A \mathbf{x} = \mathbf{b}\) com
  \[
    \pause
    A =
    \begin{pmatrix}
      2 & 1 \\
      1 & 1
    \end{pmatrix}
    \pause
    \qquad\qquad
    \mathbf{x} = \begin{pmatrix} x \\ y \end{pmatrix}
    \pause
    \qquad\qquad
    \mathbf{b} =  \begin{pmatrix} 5 \\ 3 \end{pmatrix}
  \]
\end{resolution}

%------------------------------------------------------------------------------%
\begin{resolution}{Inversa de $A$}
  Do exercício anterior
  \pause
  \[
    A^{-1} =
    \begin{bmatrix}
      1  & -1 \\
      -1 & 2
    \end{bmatrix}
  \]
\end{resolution}


%------------------------------------------------------------------------------%
\begin{resolution}{Solução}
  \begin{align*}
    \onslide<+->{\mathbf{x} = A^{-1}\mathbf{b}}
    \onslide<+->{& =
    \begin{pmatrix}
      1  & -1 \\
      -1 & 2
    \end{pmatrix}
    \begin{pmatrix}
      5\\
      3
    \end{pmatrix}}
    \\
    \onslide<+->{& =
    \begin{pmatrix}
      1\times 5 + (-1)3  \\
      (-1)5 + 2\times 3
    \end{pmatrix}}
    \\
    \onslide<+->{& =
    \begin{pmatrix}
      5-3  \\
      -5+6
    \end{pmatrix}}
    \\
    \onslide<+->{& =
    \begin{pmatrix}
      2  \\
      1
    \end{pmatrix}}
  \end{align*}

  \[
    \onslide<+->{x = 2}
    \qquad\qquad
    \onslide<+->{y = 1}
  \]
\end{resolution}

%------------------------------------------------------------------------------%
\endinput

%------------------------------------------------------------------------------%
\begin{question}
\end{question}

%------------------------------------------------------------------------------%
\begin{resolution}{Solução}
\end{resolution}

%------------------------------------------------------------------------------%
\endinput


%------------------------------------------------------------------------------%
\begin{frame}{Propriedades da matriz inversa}
  \onslide<+->{Se $A$ e $B$ são matrizes invertíveis de mesma ordem}
  \begin{gather*}
    \onslide<+->{\left(AB\right)^{-1} \;=\; B^{-1}A^{-1}}
    \\
    \onslide<+->{\left(A^{-1}\right)^{-1} \;=\; A}
    \\
    \onslide<+->{I^{-1} \;=\; I}
    \\
    \onslide<+->{\left(\lambda A\right)^{-1} \;=\; \frac{1}{\lambda} A^{-1}
    \qquad\qquad
    \text{ para um escalar } \lambda \neq 0}
    \\
    \onslide<+->{\left(A^t\right)^{-1} \;=\; \left(A^{-1}\right)^t}
  \end{gather*}
\end{frame}

%------------------------------------------------------------------------------%
\begin{frame}{Matriz Singular}
  Uma matriz quadrada $A$ que não possui inversa
  \pause
  é \emph{singular}

  \pause
  Nesse caso
  \[
    \posto(A) < n
  \]

  \pause
  \(
    A\mathbf{x} = \mnula
  \)
  possui uma solução não trivial
\end{frame}

%------------------------------------------------------------------------------%
%\section{Lista Mínima}
%
%%------------------------------------------------------------------------------%
%\begin{frame}{Lista Mínima}
%  \begin{enumerate}
%    \item
%  \end{enumerate}
%
%  \vfill
%  Atenção: A prova é baseada no livro, não nas apresentações
%\end{frame}

%------------------------------------------------------------------------------%
\end{document}
%------------------------------------------------------------------------------%


